% Avsnitt 12.4, ur lectures/L12/appendix/c_exercises.md.

\section{Övningar}
\label{c12:sec:exercises}

\begin{aside}
Övning~\ref{c12:ex:metaprev} befäster \secref{c12:sec:metaprev}; Övning~\ref{c12:ex:bittimer} till
\ref{c12:ex:drift} befäster \secref{c12:sec:bittimer}. Konstruera varje modul utifrån
specifikationen i sitt eget avsnitt; de utdelade testbänkarna, körda enligt bilaga~\ref{app:sim}, är
kontrollen av att de beter sig korrekt.
\end{aside}

% ------------------------------------------------------------------------------------------
\exerciseset{Synkroniseraren}

\exercise{\code{meta_prev}}{VHDL}
\label{c12:ex:metaprev}
Du har byggt den här modulen förut, så det mesta här handlar om de två sakerna som är nya.

\xpart{a} Skriv din egen \file{meta_prev.vhd} utifrån \secref{c12:sec:metaprev} och kör dess
testbänk:

\begin{shell}
cd controller
ghdl -a --std=93 meta_prev.vhd meta_prev_tb.vhd
ghdl -e --std=93 meta_prev_tb
ghdl -r --std=93 meta_prev_tb --assert-level=error
\end{shell}

Lägg märke till att ingen \file{can_def.vhd} står med på den första raden. \code{meta_prev} är den
enda modulen i boken som inte läser något paket alls, vilket också är skälet till att det är den som
återanvänds tre gånger.

\xpart{b} Ta bort den ena av de två vipporna, så att ett enda register blir kvar, och kör om. Vilken
kontroll fallerar, och vad säger dess meddelande dig? Sätt tillbaka den. (En version med en vippa är
inte bara svagare; den är precis den bugg testbänken finns till för att fånga, och den skulle annars
passera varje funktionell simulering i boken, eftersom ingenting i en simulator någonsin är
metastabilt.)

\xpart{c} \code{meta_prev} har ingen reset. Säg vad \code{sync_out} läser på den första klockflanken
efter uppstart, varför det är ofarligt här, och vad som skulle gå fel i \code{can_controller} självt
om modulen \emph{hade} haft en ingång \code{reset_s2_n}: instansen \code{reset_sync} är den som
producerar \code{reset_s2_n}, så vad skulle den ta för sin egen reset?

\xpart{d} \code{can_controller} instansierar den två gånger, båda med standardbredden. Säg vilka två
av dess femton portar det gäller, och varför just de två och inga andra. Säg sedan vad den tredje
instansen i designen synkroniserar \textemdash{} den som \code{can_spi_node} äger på sin toppnivå i
\chapref{c19:ch} \textemdash{} och varför just den inte kan ligga inuti \code{can_controller} i
stället.

\xpart{e} Två omkopplare slås om i samma ögonblick in i en instans med \code{WIDTH = 10}. Förklara
varför \code{sync_out} under ett kort ögonblick kan visa den ena omkopplarens nya värde och den
andras gamla, varför det inte är en bugg i \code{meta_prev}, och varför det skulle spela betydligt
större roll om de tio bitarna vore en räknare som lästes över en klockdomängräns i stället för
omkopplare som ställs för hand.

% ------------------------------------------------------------------------------------------
\exerciseset{Bittimern}

\exercise{\code{bit_timer}}{VHDL}
\label{c12:ex:bittimer}
Skriv din egen \file{bit_timer.vhd} utifrån \secref{c12:sec:bittimer}. Verifiera den sedan:

\begin{shell}
cd controller
ghdl -a --std=93 can_def.vhd bit_timer.vhd bit_timer_tb.vhd
ghdl -e --std=93 bit_timer_tb
ghdl -r --std=93 bit_timer_tb --assert-level=error
\end{shell}

Om en kontroll fallerar, läs assertion-meddelandet (bilaga~\ref{app:sim}); det namnger exakt vilket
tick och vilket förväntat värde det gäller.

\exercise{En konfigurerbar omsynkroniseringsoffset}{Simulering}
\label{c12:ex:offset}
Riktig CAN-hårdvara behöver ibland en bittimer som kan omsynkronisera \emph{och} omedelbart hoppa
fram ett fast antal tick (en förenklad ersättare för de justeringar av fassegment som riktig
CAN-bittajmingslogik utför, nämnda i \secref{c10:sec:bittiming}).

\xpart{a} Lägg till en ny generic, \code{RESYNC_OFFSET: natural := 0}, i \code{bit_timer}. När
\code{resync = '1'} ska räknaren sättas till \code{RESYNC_OFFSET} (i stället för alltid \code{0}) på
nästa klockflank.

\xpart{b} Med \code{RESYNC_OFFSET = 0} måste din ändrade entitet bete sig exakt som originalet;
bekräfta det genom att köra den \emph{oförändrade} \file{bit_timer_tb.vhd} mot den; den ska
fortfarande passera utan några ändringar i testbänken.

\xpart{c} Räkna ut på papper var \textbf{båda} utgångarna hamnar när \code{RESYNC_OFFSET} är
\code{10}: hur många tick efter en \code{resync}-puls som \code{sample} slår till, och hur många
tick efter den som \code{bit_done} gör det. Båda flyttar sig, och av samma skäl \textemdash{}
räknaren startar om på \code{10} i stället för \code{0}, så hela perioden är tio tick kortare.

Låt sedan den utdelade testbänken mäta upp det åt dig. Ändra tillfälligt din generics standardvärde
till \code{10} och kör om \code{bit_timer_tb}. Den binder ingen generic map, så den plockar upp
vilket standardvärde du än sätter, och fall 4 är det som märker det:

\begin{console}
bit_timer_tb.vhd:132:17:@6832ns:(assertion error): bit_timer: unexpected sample at tick 25 after resync!
\end{console}

Lägg märke till \textbf{vilken} av de två kontrollerna som fallerar först. Fall 4 kontrollerar inte
bara att \code{bit_done} kommer på rätt tick utan också att sampelpunkten flyttar med om, och
sampelpunkten ligger först i perioden \textemdash{} så det är den som slår till, fjorton tick innan
\code{bit_done}-kontrollen skulle ha gjort det. Under \code{--assert-level=error} stannar
simuleringen där, så meddelandet om \code{bit_done} får du aldrig se.

Meddelandet namnger själv ticket: \code{sample} kom på tick 25, alltså
\code{SAMPLE_TICK - RESYNC_OFFSET}. Vill du se \code{bit_done} också, sänk sampelkontrollens
\code{severity error} till \code{severity note} i din egen kopia av testbänken och kör om; den
rapporterar då tick 39, alltså \code{TICKS_PER_BIT - 1 - RESYNC_OFFSET}. Bekräfta att båda stämmer
med ditt pappersvar, sätt sedan tillbaka standardvärdet till \code{0} och kontrollera att den
\emph{oförändrade} testbänken passerar igen.

Det är värt att lägga märke till i sig: en testbänk som fallerar \emph{användbart} säger inte bara
att en design är fel utan också hur mycket, vilket är skälet till att assertion-meddelandena i det
här projektet alla namnger det tick eller det värde de förväntade sig.

\exercise{Hur länge kan en bittimer frilöpa?}{Resonemang}
\label{c12:ex:drift}
\code{bit_timer} omsynkroniserar \textbf{en gång per ram}, vid SOF, och frilöper därifrån
(\secref{c10:sec:stuffing}:s resonemang om klockåtervinning). Det fungerar bara om nodens egen
oscillator håller sig tillräckligt nära varenda annan nods under hela ramen. Den här övningen sätter
en siffra på ”tillräckligt nära”.

Anta att en nods 50 MHz-oscillator i själva verket går 1~\% för fort, så att dess tick är 1~\%
kortare än sändarens, och att allt annat är som specificerat: \code{TICKS_PER_BIT = 50},
\code{SAMPLE_TICK = 35}.

\xpart{a} Hur många tick har den här nodens räknare vunnit på sändarens efter en bitperiod? (Uttryck
det som en andel av ett tick.)

\xpart{b} Noden samplar vid tick 35 av 50, så i förhållande till den sanna biten kan den driva 15
tick \emph{senare} innan den samplar bortom bitens slut, eller 35 tick \emph{tidigare} innan den
samplar innan biten börjar. En snabb klocka driver sampelpunkten tidigare. Efter hur många
bitperioder har den drivit hela 35 tick?

\xpart{c} En standardram med 8 databytes går på ungefär 110--130 bitar på ledningen när stoppbitarna
är inräknade. Jämför det med ditt svar på \textbf{b)}: samplar den här noden fortfarande rätt vid
EOF? Vid ungefär vilket fält börjar den sampla fel bit?

\xpart{d} Riktiga CAN-kontrollrar omsynkroniserar \emph{under} en ram, inte bara vid SOF, med hjälp
av de fassegment och den synchronization jump width som \secref{c10:sec:bittiming} nämner. Förklara
med en eller två meningar vad det ger dem, uttryckt i ditt svar på \textbf{b)}.
